\documentclass[12pt]{article}
\usepackage[utf8]{inputenc}
\usepackage[spanish]{babel}
\usepackage{amsmath}
\usepackage{amsfonts}
\usepackage{amssymb}
\usepackage{geometry}
\usepackage{booktabs}
\usepackage{float}

\geometry{margin=1in}

\title{Reporte Ejecutivo: Plan de Lanzamiento ``Cumbia Algorítmica''}
\author{Miguel Camargo, Nicolas Cardenas y Camilo Hernandez \\ Universidad Externado de Colombia}
\date{\today}

\begin{document}

\maketitle

\section{Introducción}
Este reporte presenta la estrategia de lanzamiento para la obra experimental del artista MA, basada en un análisis estadístico bayesiano de los datos de mercado. Para garantizar objetividad, hemos priorizado el \textbf{escenario no informado}, que permite que los datos reales de los oyentes y del streaming dicten la hoja de ruta sin sesgos previos.

\section{¿La canción tiene suficiente aceptación?}
Sí, la canción es viable para un lanzamiento comercial de nicho. El análisis muestra una tasa de aceptación de aproximadamente el \textbf{51.5\%} en el escenario no informado. 
\begin{itemize}
    \item Aunque es una propuesta experimental y ``polarizante'' (un grupo considerable prefiere no escucharla), la mayoría de los oyentes potenciales coinciden en que tiene lugar en sus playlists.
    \item Esta cifra es lo suficientemente sólida para justificar una inversión inicial en promoción digital.
\end{itemize}

\section{¿En qué plataformas tiene mayor potencial?}
La plataforma con el desempeño más destacado es, sin duda, \textbf{Spotify}. Basándonos en los datos reales de 8 días de seguimiento, los resultados son contundentes:
\begin{enumerate}
    \item \textbf{Spotify}: Lidera con una media aproximada de \textbf{~5,187} reproducciones diarias .
    \item \textbf{YouTube}: Se posiciona como el segundo canal con aproximadamente \textbf{~4,245} reproducciones diarias promedio.
    \item \textbf{Apple Music}: Registra un volumen menor de aproximadamente \textbf{~2,216} reproducciones diarias promedio.
\end{enumerate}
Recomendamos concentrar el 70\% del presupuesto publicitario en Spotify, utilizando YouTube como canal de apoyo visual para reforzar el descubrimiento de la obra.

\section{¿Rock al Parque o Estéreo Picnic?}
Nuestra recomendación estratégica es el lanzamiento en \textbf{Rock al Parque}. 
\begin{itemize}
    \item Los datos muestran que el éxito de este género varía mucho día a día, lo que es típico de propuestas innovadoras que no siguen una fórmula comercial fija.
    \item Rock al Parque ofrece un público más crítico y habituado a sonidos no convencionales, lo que se alinea mejor con el perfil de la canción que un festival masivo y puramente comercial como el Estéreo Picnic.
\end{itemize}

\section{Conclusión Final}
El proyecto es viable. La combinación de una aceptación superior al 50\% y un canal de tracción fuerte en Spotify nos da luz verde para proceder con el lanzamiento. La clave del éxito será enfocar la comunicación en audiencias alternativas y melómanos que frecuentan espacios como Rock al Parque.

\end{document}
